% !TeX root = ../MAIN.tex
% !TeX encoding = UTF-8

\chapter{Einführung}
\label{chap:introduction}

Der inhaltliche Aufbau von Bachelor- bzw. Masterarbeiten, die am DKE-Institut durchgeführt werden, hat sich an folgender Aufzählung zu orientieren:

\begin{itemize}
	\item Kurzfassung (deutsch) bzw. Abstract (englisch)
	\item Einleitung mit Definition der Problemstellung (gemäß Exposé)
	\item Grundlagen / State-of-the-Art
	\item Beschreibung der Eigenleistung
	\item Fazit / Schlussfolgerungen
	\item (Anhang: ggf. Installationsanleitung, Benutzerhandbuch, etc.)
\end{itemize}

Die Beschreibung der Eigenleistung umfasst bei Implementierungsthemen folgende Punkte:

\begin{itemize}
	\item Anforderungen/Konzepte (Systemanalyse – WAS war zu implementieren?)
	\item Design bzw. Architektur der Software (Systembeschreibung/Spezifikation – WIE wurde das Geforderte umgesetzt?)
		\begin{itemize}
			\item Lösungsidee(n)
			\item Datenmodell der Software
			\item Komponenten / Module sowie deren interne und externe Abhängigkeiten
			\item Prozesse bzw. Abläufe von Algorithmen
		\end{itemize}
	\item Lösungsansatz, ggf. mit Darstellung unterschiedlicher Alternativen
	\item Schnittstelle der Software
	\item Test-/Erfahrungsbericht (z.B. Performancetests)
\end{itemize}

Die Beschreibung der Eigenleistung umfasst bei Literaturthemen folgende Punkte:

\begin{itemize}
	\item Vergleichskriterien (Kriterienkatalog, mit Motivation bzw. Begründung)
	\item Auswahl relevanter Literatur (mit Begründung)
	\item Analyse/Beschreibung der ausgewählten Ansätze in der relevanten Literatur
	\item Evaluierung der ausgewählten Ansätze anhand der Vergleichskriterien
\end{itemize}